\chapter{Summary and conclusions}\label{cap:conclusiones}

Coherent Elastic Neutrino-Nucleus Scattering (CE$\nu$NS) exhibits a cross-section scaling behavior directly proportional to the square of the weak charge ($\mathcal{Q}_W$) and, by extension, to the square of the neutron number (N) in the target material. This characteristic distinguishes CE$\nu$NS from typical neutrino interactions and highlights its significance in neutrino research. In this work, we presented an experimental strategy that used a detector array of distinct isotopes of the same element. We assumed that these detectors would operate simultaneously to obtain correlated systematic errors.
Due to their use in previous experiments, we focused on Zinc and Silicon detectors, specifically isotopes $^{64}Zn$ and $^{68}Zn$ for Zinc and $^{28}Si$ and $^{30}Si$ for Silicon. The analysis considered two distinct neutrino flux sources: one resembling an $\pi$-DAR source similar to the SNS and another one derived from a reactor flux, each with their respective characteristics. 


The differential event rate was obtained for each of the considered detectors, and it was possible to observe a specific point in the energy spectrum where the differential rates of different isotopes intersect. That reveals specific energy intervals where detectors with fewer neutrons in the target material may actually yield more CE$\nu$NS events despite the expected quadratic $N^2$ dependence. 

We explored the possibility of conducting a $\chi^2$ test to assess detector sensitivity in  CE$\nu$NS experiments regarding the quadratic dependence of the cross-section on the number of neutrons $\left(N^2\right)$. Our analysis considered the theoretical and experimental expected number of events and took into account statistical and systematic uncertainties through a covariance matrix. The $\chi^2$ analysis focused first on the $\pi$-DAR source scenario for Zinc and Silicon detectors exposed to the source for both 1-year and 5-year experiment duration. We obtained future expected allowed regions for N' consistent with the expected $N^2$ dependence. These regions form ellipses, with statistical errors dominating the width of the ellipses, while systematic errors determine their length. \\
Additionally, we extended the analysis to the reactor antineutrino source, considering different energy ranges. Figures \ref{N2EventsZn} and \ref{N2EventsSi} illustrate the allowed values for N' in both correlated and uncorrelated detector scenarios, highlighting the impact of systematic errors. The flux coming from the reactor produces a large number of events, which translates into a significant decrease in statistical error compared to the $\pi$-DAR case. This substantial statistical data significantly enhances precision. 

As an additional step, we investigated the relationship between the expected number of events and the number of neutrons in the nucleus. We showed regions where the expected permissible values of the number of events are evident for different recoil energy ranges, providing insights into whether detectors can validate the quadratic dependence experimentally.
Our analysis confirmed that the quadratic dependence has better chances of being probed at low recoil energies and suggested that Silicon is a more practical material for validating this dependence over a broader energy range.

Finally, we explored neutrino Non-Standard Interactions (NSIs) in our CE$\nu$NS proposed experimental setup.
These NSIs introduce the possibility of alternative boson mediators or other mechanisms through which neutrinos interact with matter, potentially affecting their behavior. The mathematical representation of NSIs involves parameters that quantify their strength relative to the usual weak force. These parameters can modify the weak nuclear charge and, consequently, the effective cross-section of CE$\nu$NS.  We performed an analysis based on a flux from a $\pi$-DAR source to constrain the NSI parameters. It included simultaneous constraints on non-universal parameters for electron and muon neutrino components and flavor-changing parameters. We also performed an analysis using an electron antineutrino flux from a reactor. The reactor analysis focused on constraining NSI parameters related to electron antineutrinos.

Our NSI analysis showed that, although NSI parameters exhibit degeneracies, the correlation between different isotope detectors can help to break them for the reactor antineutrino flux case, besides improving the sensitivity to these NSI parameters. 
We analyzed in detail how improving systematic uncertainties leads to tighter sensitivity. We ended our computations by presenting the combined analysis of free values for the weak mixing angle and NSI parameters. In this case, we showed that the correlation between detectors would break the corresponding degeneracies.

In summary, we have thoroughly studied the potential of an experimental
setup made of two simultaneous detectors of different isotopes of
the same material, either Zinc or Silicon. We showed their
capabilities for having a precise measurement that would allow a
better determination of standard (weak mixing angle or number of
neutrons) or beyond the Standard Model physics (NSI). We have also
made a detailed study of the perspectives in discriminating the
typical ambiguities that appear when constraining different parameters
simultaneously. We also showed that this experimental array may
disentangle these degeneracies. Our analysis is still far from an
actual experimental setup, but it is helpful as a first approximation,
and a more detailed computation taking into account feasible detectors
and neutrino sources can be pursued as a perspective of this thesis.





